\section{Fazit}
In diesem Versuch wurden die Gammastrahlen-Energiespektren der drei untersuchten Elemente ($^{152}$Eu, $^{60}$Co und $^{137}$Cs) sowie von Gesteinsproben, die  mitgebracht wurden, aufgezeichnet und verglichen. Die Messungen erfolgten mit zwei verschiedenen Detektoren: einem NaI(Tl)-basierten Szintillationsdetektor (Sz-Detektor) und einem hochreinen Germaniumdetektor (HPGe-Detektor).

Für die Energiekallibration wurden die Thermschemata der radioaktiven Isotope verwendet. Bei beiden Detektoren wurde ein linearer Zusammenhang zwischen Kanalnummer und Photonenenergie beobachtet, was auf einen linearen Kalibrierungsprozess hindeutet.

Anschließend wurde die Halbwertsbreite (FWHM) der Spektrallinien analysiert, um die Signalempfindlichkeit zu charakterisieren. (Für den HPGe-Detektor wurde eine Abhängigkeit der FWHM beobachtet, die mit den erwarteten Ergebnissen der Ladungsträgerfunktionalisierung übereinstimmt. )

Das Effizienzverhältnis wurde ebenfalls bestimmt. Es erwies sich als ????? wie das des Sz-Detektors. Dies lässt sich auf Unterschiede in der Ordnungszahl, dem Photonenfluss und der Geometrie der aktiven Teilchen zurückführen/ weiter Kommentare.

Darüber hinaus wurde die Leistungsfähigkeit der beiden Verbindungen ermittelt. Der HPGe-Detektor liefert eine Leistung von
\begin{align*}
    \mu = (\pm)\%.
\end{align*}
Der Sz-Detektor weist aufgrund der größeren Größe des NaI(Tl)-Kristalls und seiner höheren effektiven Ordnungszahl eine etwa ???-fach höhere Empfindlichkeit auf. Die Effizienzkurve wurde mithilfe der relativen Effizienzmethode erstellt. Die resultierenden Signale lassen sich über die Anregungsenergie anpassen, und ihre Morphologie stimmt mit der Energiebilanz des photoelektrischen und des Compton-Effekts überein. Daher repräsentiert das Verhältnis das erwartete Energieniveau für eine optimale Leistung.

Die Effizienzkurve wurde verwendet, um die Aktivität der beiden $^{60}$Co-Peaks zu bestimmen. Folgende Ergebnisse wurden erzielt:

???? keV-Linie:
\begin{align*}
 A = (\pm) \tfrac{1}{s}   
\end{align*}l


???? keV-Linie:

\begin{align*}
 A = (\pm) \tfrac{1}{s}   
\end{align*}l
Kommentare zu den Ergebnissen.

Abschließend wurde das Gammaspektrum der Gesteinsproben analysiert. Die Intensität der beobachteten Spektrallinie kann durch Kalibrierung präzise bestimmt werden. Ein Vergleich mit der Literatur zeigt eine gute Übereinstimmung mit den Gammalinien der folgenden Kerne: (235A, 227B, 228Th, 224Ra, 210Rn). Darüber hinaus lässt sich aus den beobachteten Zerfallsketten auf das Vorhandensein von (232Th) in der Probe schließen.